\documentclass[10pt,a4paper]{article}
\usepackage[spanish,es-nodecimaldot]{babel}
\usepackage[utf8]{inputenc}
\usepackage[T1]{fontenc}
\usepackage{lmodern}
\usepackage[top=1.1cm,bottom=1.1cm,left=1.5cm,right=1.5cm]{geometry}
\usepackage{amsmath,amssymb}
\usepackage[table]{xcolor}
\usepackage{enumitem}
\usepackage[normalem]{ulem}
\usepackage{tikz}
\usetikzlibrary{arrows.meta}

\setlength{\parindent}{0pt}
\setlength{\parskip}{2pt}
\pagestyle{empty}

\AtBeginDocument{%
  \setlength{\abovedisplayskip}{3pt}%
  \setlength{\belowdisplayskip}{3pt}%
  \setlength{\abovedisplayshortskip}{2pt}%
  \setlength{\belowdisplayshortskip}{2pt}%
}

\begin{document}

\begin{center}
{\large\bfseries Inferencias basadas en min cuadrados}
\end{center}
\vspace{-2pt}

Una vez que calculamos a y b para nuestra muestra, la inferencia basada en min cuadrados nos permetiran saber que tan confiable es ese modelo para toda la poblacion y no para una muestra en particular. La inferencia sirve para pasar de los estimadores muestrales \textbf{a} y \textbf{b} a conclusiones sobre los parametros poblacionales $\boldsymbol{\alpha}$ y $\boldsymbol{\beta}$.

\textbf{Sxy , Sxx} me van a permitir estimar la pendiente b y la ordenada al origen a. Con \textbf{Syy} voy a poder calcular la varianza residual muestral

\begin{center}
{\small (1)} \textbf{Calc.--> Sxy, Sxx, Syy}
\qquad\qquad
{\small (2)} $\displaystyle \bar{y} = \frac{\sum y_i}{n} \qquad \bar{x} = \frac{\sum x_i}{n}$
\end{center}

\begin{center}
{\small (3)} \quad $\displaystyle b = \frac{S_{xy}}{S_{xx}} \qquad\qquad a = \bar{Y} - b\bar{X}$
\end{center}

\begin{minipage}[t]{0.46\textwidth}
\vspace{0pt}
{\small (4)}
\[
s_e^2 = \frac{S_{xx}\cdot S_{yy} - S_{xy}^2}{n(n-2)\cdot S_{xx}}
\qquad \text{\footnotesize n-2 gl}
\]
\centerline{Luego}
\[
s_e = \sqrt{s_e^2}
\]
\end{minipage}\hfill
\begin{minipage}[t]{0.50\textwidth}
\vspace{6pt}
{\small
$s_e^2 \longrightarrow$ Varianza residual muestral

$s_e \longrightarrow$ Error estandar de estimacion

\textbf{n-->cant. puntos}
}
\end{minipage}

\centerline{Intervalos de conf. $\alpha$ y $\beta$}

\begin{center}
\begin{minipage}{0.42\textwidth}
\centering
$\alpha$
\[
t = \frac{a-\alpha}{s_e}\cdot \sqrt{\frac{n\cdot S_{xx}}{S_{xx}+(n\cdot\bar{x})^2}}
\]
\end{minipage}
\begin{minipage}{0.42\textwidth}
\centering
$\beta$
\[
t = \frac{b-\beta}{s_e}\cdot \sqrt{\frac{S_{xx}}{n}}
\]
\end{minipage}
\hfill {\small $\longrightarrow$ Se compara con el t.crit}
\end{center}
\vspace{-4pt}
\hfill{\small t. calc $\uparrow$}\hspace{5.2cm}

\[
a - t_{\frac{\alpha}{2}}\cdot s_e\sqrt{\frac{n\cdot S_{xx}}{S_{xx}+(n\cdot\bar{x})^2}}
\;<\; \alpha \;<\;
a + t_{\frac{\alpha}{2}}\cdot s_e\sqrt{\frac{n\cdot S_{xx}}{S_{xx}+(n\cdot\bar{x})^2}}
\]
\[
b - t_{\frac{\alpha}{2}}\cdot s_e\sqrt{\frac{n}{S_{xx}}}
\;<\; \beta \;<\;
b + t_{\frac{\alpha}{2}}\cdot s_e\sqrt{\frac{n}{S_{xx}}}
\qquad
{\footnotesize\textcolor{red!70!black}{\text{Errores máximos de estimación de a y b}}}
\]

Las pruebas de hipótesis con $\boldsymbol{\beta}$ son muy importantes, pq indica el cambio promedio de \textbf{Y} cuando \textbf{X} aumenta 1 unidad. Si $\beta=0$: y no depende de x (recta horizontal), si $\beta \neq 0$ y depende de x por lo que existe relacion lineal.

Habiendo comprobado q $\beta \neq 0$ osea que existe relacion lineal, podemos utilizar dicha recta ajustada para realizar estimaciones. Podemos ver q pasa en un punto especifico X0

\vspace{4pt}
\hrule
\vspace{4pt}

\begin{center}
{\large\bfseries Inferencias basadas en min. cuadrados}
\end{center}
\vspace{-2pt}

Una vez que calculamos a y b para nuestra muestra, la inferencia basada en min cuadrados nos permetiran saber que tan confiable es ese modelo para toda la poblacion y no para una muestra en particular. La inferencia sirve para pasar de los estimadores muestrales \textbf{a} y \textbf{b} a conclusiones sobre los parametros poblacionales $\boldsymbol{\alpha}$ y $\boldsymbol{\beta}$. Nos va a permitir :

\textbf{1)Probar hipotesis sobre la pendiente $\boldsymbol{\beta}$ (la + importante!)}

\begin{minipage}[t]{0.24\textwidth}
\vspace{0pt}
\[
H_0 : \beta = 0
\]
\[
H_1 : \beta \neq 0
\]
\end{minipage}
\begin{minipage}[t]{0.72\textwidth}
\vspace{0pt}
{\small
*Si no rechazamos H0, no hay evidencia suficiente de una relación lineal entre X e Y.

*Si rechazamos H0, concluimos que existe evidencia de que X influye linealmente sobre Y.

La prueba sobre $\beta$ es la más importante porque $\beta$ indica cuánto cambia la respuesta media Y cuando X aumenta una unidad.
}
\end{minipage}

\textbf{2)Probar hipótesis sobre la ordenada al origen $\boldsymbol{\alpha}$ (no importa mucho}, para saber si la recta poblacional pasa o no por el origen)

\textbf{3)Construir intervalos de confianza para $\boldsymbol{\alpha}$ y $\boldsymbol{\beta}$.} Lo que nos permite afirmar entre que valores probablemente se ecuentra la verdadera pendiente poblacional $\beta$

\textbf{4) Estimar la respuesta media para un valor especifico x0}

\hspace{1cm}$\hat{y}_0 = a + bx_0$ \quad Luego se construye un intervalo de confianza para la respuesta media poblacional cuando X=x0.

\end{document}
